% Part: applied-modal-logics
% Chapter: epistemic-logic
% Section: properties-accessibility

\documentclass[../../../include/open-logic-section]{subfiles}

\begin{document}

\olfileid{aml}{el}{acc}

\olsection{رسائی دے تعلق تے علم دے اصول}

معیاری امکانی منطقاں وچ فریم دی خاصیت تے فارمولے دے میل بارے
اسیں پہلاں ای کجھ جاندے آں۔ ہُن ویکھ سکدے آں کہ علم دے لحاظ نال
تعبیر کرن تے خاصیت دی پچھان دین والے !!{formula}s کِہو جیہے ہوندے
نیں۔ اسیں آکھ چکے آں کہ علم دیاں منطقاں دی تعبیر عام طور اُتے ایس 5
وچ ہوندی اے۔ سو ویکھدے آں کہ علم دے وکھ وکھ اصول کِویں وکھائے
جاندے نیں تے فریم دیاں کیہڑیاں شرطاں نال ملدے نیں۔

معیاری امکانی منطق توں یاد کرو کہ وکھ وکھ امکانی !!{formula}s
رسائی تعلقاں دیاں وکھ وکھ خاصیتاں دی پچھان دیندے سن۔ ایس جدول وچ
خاص علم اصولاں نال ملدیاں کجھ خاصیتاں چنیاں نیں۔

\begin{table}[t]
    \begin{tabular}{| p{.48\textwidth} || p{.48\textwidth} |}
      \hline
      {\emph{جے $R$ \dots اے}} & {\emph{تاں \dots،~$\mModel{M}$ وچ سچ اے}} \\
      \hline \hline

      & $\Knows (p \lif q) \lif (\Knows p \lif \Knows q)$ 
      \hfill \newline{(استلزام تحت بندش)} \\
      \hline
      \emph{انعکاسی}: $\forall w Rww$  
      & $\Knows p \lif p$ \hfill \newline{(علم دی صداقت)} \\
      \hline
      \emph{متعدی}: \newline
      $\forall u \forall v \forall w ((Ruv \land Rvw) \lif Ruw)$ & 
      $\Knows p \lif \Knows \Knows p$ \hfill 
           \newline{(مثبت خود آگاہی)} \\
      \hline 
      \emph{اقلیدسی}: \newline
      $\forall w \forall u \forall v ((Rwu \land Rwv) \lif Ruv)$ &  
      $\lnot \Knows p \lif \Knows \neg \Knows p$ \hfill 
        \newline{(منفی خود آگاہی)} \\
      \hline
    \end{tabular}
    \caption{علم دے چار اصول۔}
    \ollabel{tab:four}
  \end{table} 

علم دی صداقت، جیہڑی مسلّمے $T$ نال ملدی اے، اکثر ایہناں اصولاں
وچوں سبھ توں گھٹ اختلاف والی سمجھی جاندی اے۔ ایہ آکھدی اے کہ جے
!!a{formula} دا پتا اے، تاں اوہ لازمی سچ اے۔ استلزام تحت بندش تے
مثبت تے منفی خود آگاہی اُتے کِتے ودھ اختلاف اے۔

استلزام تحت بندش، جیہڑی مسلّمے $K$ نال ملدی اے، ایہ خیال وکھاؤندی
اے کہ فاعل دا علم استلزام تحت بند اے۔ کجھ صورتاں وچ ایہ گل منّن
جوگی لگ سکدی اے۔ مثلاً مینوں پتا ہو سکدا اے کہ جے میں وکٹوریا وچ
آں، تاں وینکوور ٹاپو اُتے آں۔ جے شک دیاں انوکھی صورتاں پاسے رکھیے،
تاں مینوں پتا اے کہ میں وکٹوریا وچ آں۔ ایس توں ایہ وی سمجھ آؤندی اے
کہ مینوں پتا اے کہ میں وینکوور ٹاپو اُتے آں۔ سو ایس صورت وچ میرے
علم دی منطقی بندش چنگی طرح سمجھ آ سکدی اے۔ دوجے پاسے اسیں اپنے علم
دے نتیجیاں اُتے ہمیشہ پورا غور نہیں کردے۔ ایس کرکے ہور صورتاں وچ
ایسے نتیجے نکل سکدے نیں جیہڑے گھٹ سِدھے لگن۔

مثبت خود آگاہی نوں کدی کدی کے کے دا اصول آکھیا جاندا اے۔
اینوں ایویں بیان کیتا جاندا اے کہ جے مینوں کسے گل دا پتا اے، تاں
مینوں ایہ وی پتا اے کہ مینوں اوہدا پتا اے۔ ایہ مسلّمے 4 دا علم
والا روپ اے۔ ایس طرح منفی خود آگاہی آکھدی اے کہ جے مینوں کسے
گل دا پتا \emph{نہیں}، تاں مینوں پتا اے کہ مینوں اوہدا پتا نہیں۔
ایہ مسلّمے 5 دا علم والا روپ اے۔ دونہاں دیاں عام جہیاں خلاف مثالاں
لگدیاں نیں، جتھے مینوں یقین نہیں کہ مینوں کسے گل دا پتا اے یا
نہیں، بھاویں اصل وچ مینوں اوہدا پتا اے۔

\end{document}
